\section{Conclusão}
Tendo como base o conhecimento desenvolvido neste relatório, pode-se concluir que ferramentas como a transformada Z e sua inversa são de grande utilidade na análise de sistemas complexos, tendo em vista que elas, na maioria dos casos, simplificam as expressões algébricas e permitem uma análise mais facilitada do sistema.

Outrossim, para verificar como o sistema se comporta no início de sua operação e após um delongado período de tempo, nota-se que os teoremas do valor final e do valor inicial são de extrema importância, em especial no domínio discreto, como foi avaliado neste relatório, sendo isso validado na seção \ref{resultado}.

Além disso, a aplicação dos teoremas citados anteriormente ainda permite avaliar se um sistema é estável ou não, tendo em vista que sistemas estáveis não apresentam valores ilimitados em seu funcionamento quando aplicada uma entrada estável. Ademais, pode-se avaliar a estabilidade a partir dos polos, tendo em vista que, como mencionado em \ref{estabilidade}, se todos os polos do sistema estão contidos dentro do círculo unitário, o sistema é estável.

Por fim, pode-se afirmar que as ferramentas para a realização de cálculos e análises no domínio discreto são de grande ajuda. Em especial em laboratórios e experimentos, tendo em vista que todas as amostras coletadas são valores discretos e como visto neste relatório, os meios para calcular e analisar esses sistemas, como a transformada Z e sua inversa, são eficientes.
